\begingroup
\makeatletter\tagsleft@true\let\veqno\@@leqno\makeatother
\clearpage\SourceLeaf{322b}\ArithmeticHead{315}{R}{{\scriptsize ABBILDUNG DER ELLIPSE AUF DEN KREIS.}}
\noindent {\small[S. 222]}
\begin{center}\bfseries Gründe auf denen die Reihen S. 176, 177 beruhen.\end{center}
\NumberedItem{5}
\par\indent Für die Ellipse, deren halbe grosse und kleine Axe $1+\lambda,\ 1-\lambda$ ist, sei der complexe Ausdruck eines unbestimmten Punkts
\begin{equation}
(1+\lambda)\cos e+(1-\lambda)\sin e.\,i=u+\lambda u^{-1}.\tag*{[8]}
\end{equation}
\noindent Für einen zweiten Punkt seien $E,U$ dasselbe, was $e,u$ für jenen. Der Logarithme der Entfernung ist also der reelle Theil von
\[
\log(U-u)+\log\left(1-\frac{\lambda u^{-1}}{U}\right)
\]
\noindent oder von
\[\begin{aligned}
&\frac uU-\frac12\frac{uu}{UU}-\frac13\frac{u^3}{U^3}\,{.}{.}{.}\\
&-\frac\lambda{uU}-\frac12\frac{\lambda\lambda}{uuUU}-\frac13\frac{\lambda^3}{u^3U^3}\,{.}{.}{.}\text{[**)]},
\end{aligned}\]
\noindent folglich
\begin{equation}
\left\{\begin{aligned}
&=-\cos(E-e)-\tfrac12\cos2(E-e)-\tfrac13\cos3(E-e)-\text{etc.}\\
&\quad-\lambda\cos(E+e)-\tfrac12\lambda\lambda\cos2(E+e)-\tfrac13\lambda^3\cos3(E+e)-\text{etc.}
\end{aligned}\right.\tag*{[9]}
\end{equation}
\par\indent Ist auf der Peripherie der Ellipse die Masse $2\pi A$ so vertheilt, dass auf jedes Element die Masse
\par\bigskip\noindent\rule{.25\linewidth}{.4pt}\par\smallskip
{\small\par\indent [**)] Der Kunstgriff bleibt anwendbar, wenn die zwei Punkte in verschiedenen Ellipsen mit gleichen Brennpunkten liegen.\par}
\par\medskip\noindent\hfill$40^{\ast}$
\clearpage\SourceLeaf{323}\ArithmeticHead{316}{L}{ANALYSIS.\quad NACHLASS.}
\begin{equation}
de\left\{\begin{aligned}
&A+A'\cos e+A''\cos2e+A'''\cos3e+\text{etc.}\\
&\quad+B'\sin e+B''\sin2e+B'''\sin3e+
\end{aligned}\right\}\tag*{[10]}
\end{equation}
\noindent kommt, so ist das Potential einer Masse, in dem Punkte der Ellipse, auf den sich $E$ bezieht,
{\small\begin{equation}
-\pi\left\{\begin{aligned}
&(1+\lambda)A'\cos E+\tfrac12(1+\lambda\lambda)A''\cos2E+\tfrac13(1+\lambda^3)A'''\cos3E+{.}{.}{.}\\
&+(1-\lambda)B'\sin E+\tfrac12(1-\lambda\lambda)B''\sin2E+\tfrac13(1-\lambda^3)B'''\cos3E
\end{aligned}\right\}\tag*{[11]}
\end{equation}}
\noindent Ist hingegen dieselbe Masse im Centrum concentrirt, so ist das Potential der reelle Theil von
\[2\pi A\log(U+\lambda U^{-1}),\]
\noindent [also]
\begin{equation}
=2\pi A(\lambda\cos2E-\tfrac12\lambda\lambda\cos4E+\tfrac13\lambda^3\cos6E-{.}{.}{.}).\tag*{[12]}
\end{equation}
\noindent {\small[S. 223]}
\par\indent Damit beide Potentiale gleich gross sein, muss bei der Vertheilung auf jedes Element der Peripherie der Ellipse kommen die Masse
\begin{equation}
A\,de.\left(1-\frac{4\lambda}{1+\lambda\lambda}\cos2e+\frac{4\lambda\lambda}{1+\lambda^4}\cos4e-\frac{4\lambda^3}{1+\lambda^6}\cos6e\,{.}{.}{.}\right);\tag*{[13]}
\end{equation}
\noindent wird diese Masse $=A\,dp$ gesetzt und $p$ von $e=0$ an gezählt, so ist
\[
p=e-\frac{2\lambda}{1+\lambda\lambda}\sin2e+\frac22.\frac{\lambda\lambda}{1+\lambda^4}\sin4e-\frac23.\frac{\lambda^3}{1+\lambda^6}\sin6e\,{.}{.}{.}
\]
\noindent welches mit der Formel S. 177 [der Handschrift, siehe Gleichung [7], oben S. 314] übereinstimmt, wenn $ip$ den imaginären Theil von $\log t$ bedeutet.
\clearpage
\endgroup
